\documentclass[aspectratio=169]{beamer}

% -- Theme --------------------------------------------------------------------
\usetheme{default}
\usecolortheme{default}
\setbeamertemplate{navigation symbols}{}
\setbeamertemplate{footline}[frame number]
\setbeamertemplate{itemize items}[circle]

\definecolor{mainblue}{RGB}{0,70,127}
\setbeamercolor{frametitle}{fg=mainblue}
\setbeamercolor{title}{fg=mainblue}
\setbeamercolor{itemize item}{fg=mainblue}
\setbeamercolor{itemize subitem}{fg=mainblue}
\setbeamerfont{frametitle}{size=\large, series=\bfseries}

% -- Packages -----------------------------------------------------------------
\usepackage{booktabs}
\usepackage{amsmath}
\usepackage{xcolor}
\usepackage{colortbl}

\newcommand{\fhalf}{F\textsubscript{0.5}}
\newcommand{\errtag}[1]{\texttt{#1}}

% -- Metadata -----------------------------------------------------------------
\title{Advancements in Arabic Grammatical Error\\
Detection and Correction: An Empirical Investigation}
\author{Authors: Alhafni, Inoue, Khairallah \& Habash \\ \small NYU Abu Dhabi}
\date{Paper from: EMNLP 2023}

% -----------------------------------------------------------------------------
\begin{document}
% -----------------------------------------------------------------------------

\begin{frame}
  \titlepage
\end{frame}

% -- MOTIVATION ---------------------------------------------------------------

\begin{frame}{Why Arabic GEC is Hard}
  GEC research has focused almost entirely on English.
  Arabic presents unique challenges:

  \begin{itemize}
    \item \textbf{Morphological richness:} Arabic inflects for gender, number,
      person, case, mood, voice, and aspect, producing a very large vocabulary
 \item \textbf{Two-tier language system:} Arabic has a formal written
  standard and a range of spoken dialects that constantly mix together,
  even in educated writing
\item \textbf{Built-in spelling variation:} certain letter forms are
  routinely confused even by native speakers, meaning what counts as
  an ``error'' is not always clear-cut
    \item \textbf{Punctuation:} punctuation errors make up 40\% of errors in
      Arabic GEC data -- ten times higher than in English learner corpora
    \item \textbf{Data scarcity:} existing Arabic GEC datasets have no
      error type annotations, and training sets are tiny by English standards
  \end{itemize}
\end{frame}

\begin{frame}{Three Contributions in One Paper}
  \begin{itemize}
    \item \textbf{First benchmark} of pretrained Seq2Seq Transformer models
      on Arabic GEC across three datasets\\[1em]
    \item \textbf{First multi-class Arabic GED results,} including a new
      annotation pipeline to label error types in existing datasets
      that have no type information\\[1em]
    \item \textbf{Systematic study} of whether GED improves Arabic GEC,
      including morphological preprocessing as an orthogonal intervention
  \end{itemize}
\end{frame}

% -- DATA ---------------------------------------------------------------------

\begin{frame}{Three Arabic GEC Datasets}
  \begin{center}
  \begin{tabular}{llll}
    \toprule
    \textbf{Dataset} & \textbf{Type} & \textbf{Domain} & \textbf{Training size} \\
    \midrule
    QALB-2014 & L1 (native) & News comments & 1M words \\
    QALB-2015 & L2 (learner) & Essays & 43K words \\
    ZAEBUC    & L1 (native) & University essays & 25K words \\
    \bottomrule
  \end{tabular}
  \end{center}

  \vspace{1em}
  All three have roughly 25--30\% erroneous words.
  None have manual error type annotations.
  QALB-2015 and ZAEBUC are very small by GEC standards,
  which constrains every architectural choice toward simplicity.
\end{frame}

% -- GED PIPELINE -------------------------------------------------------------

\begin{frame}{Building an Arabic GED Pipeline}
  No ERRANT equivalent exists for Arabic.
  The paper builds one from scratch in two steps.\\[1.5em]

  \textbf{Step 1: Edit extraction.}
  Align erroneous and corrected sentence pairs to locate the edits.
  Standard Levenshtein and the M2 scorer both fail on merge and split
  operations. A new iterative algorithm greedily handles multi-token edits,
  substantially outperforming all alternatives (AER 0.00 vs 0.10--0.16).\\[1em]

  \textbf{Step 2: Error type annotation.}
  Pass extracted edits to ARETA, an Arabic error type annotation tool for MSA.
  ARETA's own alignment is replaced with the new algorithm to increase coverage.\\[1em]

  Error types follow ARETA's taxonomy: orthography, morphology, syntax,
  semantics, punctuation, merges, and splits -- at three granularity levels
  (2-class, 13-class, 43-class).
\end{frame}

% -- GEC MODELS ---------------------------------------------------------------

\begin{frame}{GEC Models: Baselines}
  Four baselines are compared before introducing the main Seq2Seq models:

  \begin{itemize}
    \item \textbf{Morphological disambiguation (Morph):} CAMeL Tools
      runs a morphological analyzer and disambiguator on the input,
      producing spelling corrections, POS tags, and lemmas.
      No training data needed.
    \item \textbf{MLE:} A bigram lookup model that maps erroneous words
      to corrections using the training data. Highest precision, low recall.
    \item \textbf{ChatGPT (GPT-3.5):} 3-shot prompted. Highest recall of
      the baselines but lower precision. 38\% of mismatches are
      grammatically or lexically incorrect.
  \end{itemize}
\end{frame}

\begin{frame}{GEC Models: Seq2Seq with GED}
  Two pretrained Arabic Seq2Seq Transformers are the main models:

  \begin{itemize}
    \item \textbf{AraBART:} pretrained on 24GB of MSA data, mostly news
    \item \textbf{AraT5:} pretrained on 256GB of MSA and Twitter data
  \end{itemize}

  \vspace{1em}
  GED predictions are incorporated as a \textbf{token embedding layer}
  summed with the standard token and positional embeddings before the
  encoder's attention layers.\\[1em]

  This is simpler than Yuan et al.'s second encoder approach.
  With small training sets, not adding parameters matters:
  a second encoder would risk overfitting.\\[1em]

  Predicted (not gold) GED tags are used during training
  to match the noise present at inference time.
\end{frame}

% -- RESULTS ------------------------------------------------------------------

\begin{frame}{Does Morphological Preprocessing Help?}
  Yes, consistently. Adding Morph improves \fhalf\ across MLE, AraT5,
  and AraBART on most datasets.\\[1em]

  The gains come almost entirely from \textbf{recall}.
  Morph helps the model find more errors to fix,
  at a small cost to precision.\\[1em]

  The exception is QALB-2015-L2 where Morph has no effect or slightly
  hurts -- the L2 error profile (more morphological and syntactic errors)
  may not align well with what the morphological disambiguator corrects.
\end{frame}

\begin{frame}{Does GED Help Arabic GEC?}
  Yes. Adding GED predictions as auxiliary input improves results
  across all three datasets for both AraT5 and AraBART.
  AraBART benefits more.\\[1em]

  \begin{center}
  \begin{tabular}{lcccc}
    \toprule
    \textbf{System} & \textbf{QALB-14} & \textbf{QALB-15} & \textbf{ZAEBUC} & \textbf{Avg.} \\
    \midrule
    AraBART & 78.7 & 61.2 & 83.4 & 74.4 \\
    +Morph & 78.8 & 61.7 & 83.6 & 74.7 \\
    +GED43 & 79.1 & 61.9 & 83.9 & 75.0 \\
    \textbf{+Morph+GED43} & \textbf{79.3} & \textbf{62.4} & \textbf{84.2} & \textbf{75.3} \\
    \bottomrule
  \end{tabular}
  \end{center}

  \vspace{0.5em}
  As in English, improvements are driven by recall.
  Combining Morph and GED gives the best results on average.
\end{frame}

\begin{frame}{GED Granularity: No Clear Winner}
  Unlike Yuan et al., no single granularity level consistently wins.\\[1em]

  \begin{center}
  \begin{tabular}{lcccc}
    \toprule
    \textbf{GED granularity} & \textbf{QALB-14} & \textbf{QALB-15} & \textbf{ZAEBUC} & \textbf{Avg.} \\
    \midrule
    2-class & 79.1 & 61.9 & 82.9 & 74.6 \\
    13-class & 79.4 & 62.3 & 84.5 & 75.4 \\
    43-class & 79.3 & 62.4 & 84.2 & 75.3 \\
    \midrule
    43-class Oracle & 82.8 & 69.8 & 88.1 & 80.2 \\
    \bottomrule
  \end{tabular}
  \end{center}

  \vspace{0.5em}
  13-class edges out 43-class on average (75.4 vs 75.3).
  The large oracle gap confirms the bottleneck is GED accuracy,
  not model capacity or label granularity.
\end{frame}

\begin{frame}{State-of-the-Art Results}
  The best system achieves SOTA on two of the three benchmarks.\\[1em]

  \begin{center}
  \begin{tabular}{lcccc}
    \toprule
    \textbf{System} & \textbf{QALB-14 F1} & \textbf{QALB-15-L2 F1} \\
    \midrule
    Watson et al.\ (2018) -- prior SOTA & 70.4 & 73.2 \\
    Solyman et al.\ (2022) & 76.0 & -- \\
    \textbf{AraBART+Morph+GED (this work)} & \textbf{79.5} & \textbf{77.1} \\
    \bottomrule
  \end{tabular}
  \end{center}

  \vspace{1em}
  Results on ZAEBUC establish a strong benchmark for future work
  as no prior results exist on this dataset.
\end{frame}

% -- CONCLUSION ---------------------------------------------------------------

\begin{frame}{Summary}
  \begin{itemize}
    \item A new edit extraction algorithm handles Arabic merge and split
      operations, enabling reliable automatic error type annotation\\[0.8em]
    \item Pretrained Seq2Seq models (AraBART, AraT5) substantially
      outperform all prior Arabic GEC systems\\[0.8em]
    \item GED improves Arabic GEC in the same way it improves English GEC,
      mainly through recall gains\\[0.8em]
    \item Morphological preprocessing and GED are complementary interventions\\[0.8em]
    \item No single GED granularity wins across all datasets -- the bottleneck
      is GED accuracy, not label richness\\[0.8em]
    \item SOTA on QALB-2014 and QALB-2015; strong benchmark on ZAEBUC
  \end{itemize}
\end{frame}

\end{document}
